% figures/vaswani_transformer.tex -- Transformer-Architektur (eigene
% Darstellung in Anlehnung an Vaswani et al. 2017, TikZ).
\begin{tikzpicture}[
  font=\sffamily\footnotesize,
  blk/.style={draw=gray!60, rounded corners=2pt, minimum width=3.5cm,
    minimum height=0.64cm, align=center, inner sep=2pt},
  attn/.style={blk, fill=orange!22}, ffn/.style={blk, fill=blue!12},
  norm/.style={blk, fill=green!16, minimum height=0.5cm},
  emb/.style={blk, fill=red!10}, lin/.style={blk, fill=violet!14},
  smx/.style={blk, fill=yellow!22, minimum height=0.5cm},
  stack/.style={draw=gray!50, rounded corners=4pt, fill=gray!6, inner sep=5pt},
  arr/.style={-{Latex[length=2mm]}, thick, black!75},
  res/.style={-{Latex[length=2mm]}, thick, black!75, rounded corners=2pt},
  pe/.style={circle, draw=gray!60, inner sep=1pt, minimum size=0.42cm, fill=white},
]
  % ---------------- Encoder (x = 0)
  \node[emb]  (ein) at (0, -0.15) {Eingabe-Embedding};
  \node[attn] (esa) at (0, 1.9)  {Multi-Head-\\Self-Attention};
  \node[norm] (en1) at (0, 2.75) {Add \& Norm};
  \node[ffn]  (eff) at (0, 3.75) {Feed-Forward-\\Netz};
  \node[norm] (en2) at (0, 4.6)  {Add \& Norm};
  \begin{scope}[on background layer]
    \node[stack, fit={(esa)(en2)(-2.05, 1.3)}] (enc) {};
  \end{scope}
  \node[font=\sffamily\small, anchor=east] at ($(enc.west)+(-0.25,0)$) {$N\times$};
  \node (einp) at (0, -1.3) {Eingaben};
  \node[pe] (eplus) at (0, 0.72) {$+$};
  \node[pe] (epos) at (-2.0, 0.72) {$\sim$};
  \node[align=center, font=\sffamily\scriptsize, anchor=east] at (-2.3, 0.72) {Positions-\\kodierung};
  \draw[arr] (einp) -- (ein);
  \draw[arr] (ein) -- (eplus); \draw[arr] (epos) -- (eplus); \draw[arr] (eplus) -- (esa);
  \draw[arr] (esa) -- (en1); \draw[arr] (en1) -- (eff); \draw[arr] (eff) -- (en2);
  \draw[res] (0, 1.3)  -- ++(-2.05, 0) |- (en1.west);
  \draw[res] (0, 3.2)  -- ++(-2.05, 0) |- (en2.west);
  % ---------------- Decoder (x = 6.6)
  \node[emb]  (aus) at (6.6, -0.15) {Ausgabe-Embedding};
  \node[attn] (dma) at (6.6, 1.9)  {Maskierte Multi-Head-\\Self-Attention};
  \node[norm] (dn1) at (6.6, 2.75) {Add \& Norm};
  \node[attn] (dca) at (6.6, 3.75) {Encoder-Decoder-\\Attention};
  \node[norm] (dn2) at (6.6, 4.6)  {Add \& Norm};
  \node[ffn]  (dff) at (6.6, 5.6)  {Feed-Forward-\\Netz};
  \node[norm] (dn3) at (6.6, 6.45) {Add \& Norm};
  \begin{scope}[on background layer]
    \node[stack, fit={(dma)(dn3)(8.65, 1.3)}] (dec) {};
  \end{scope}
  \node[font=\sffamily\small, anchor=west] at ($(dec.east)+(0.25,0)$) {$N\times$};
  \node[lin]  (lin)  at (6.6, 7.55) {Lineare Schicht};
  \node[smx]  (soft) at (6.6, 8.35) {Softmax};
  \node[align=center] (outp) at (6.6, 9.3) {Ausgabe-\\wahrscheinlichkeiten};
  \node[align=center] (ausp) at (6.6, -1.3) {Ausgaben\\(rechtsverschoben)};
  \node[pe] (dplus) at (6.6, 0.72) {$+$};
  \node[pe] (dpos) at (8.6, 0.72) {$\sim$};
  \node[align=center, font=\sffamily\scriptsize, anchor=west] at (8.9, 0.72) {Positions-\\kodierung};
  \draw[arr] (ausp) -- (aus);
  \draw[arr] (aus) -- (dplus); \draw[arr] (dpos) -- (dplus); \draw[arr] (dplus) -- (dma);
  \draw[arr] (dma) -- (dn1); \draw[arr] (dn1) -- (dca); \draw[arr] (dca) -- (dn2);
  \draw[arr] (dn2) -- (dff); \draw[arr] (dff) -- (dn3); \draw[arr] (dn3) -- (lin);
  \draw[arr] (lin) -- (soft); \draw[arr] (soft) -- (outp);
  \draw[res] (6.6, 1.3)  -- ++(2.05, 0) |- (dn1.east);
  \draw[res] (6.6, 3.2)  -- ++(2.05, 0) |- (dn2.east);
  \draw[res] (6.6, 5.05) -- ++(2.05, 0) |- (dn3.east);
  % ---------------- Encoder-Ausgabe -> Encoder-Decoder-Attention (Keys, Values)
  \draw[arr] (en2.north) -- ++(0, 0.45) -| (3.3, 3.75) -- (dca.west);
  \node[font=\sffamily\scriptsize, anchor=south, inner sep=2pt] at (1.65, 5.3) {Encoder-Ausgabe (Keys, Values)};
\end{tikzpicture}
